%------------------------------------------------------------------------------%
\begin{question}
  Determine o ângulo entre a reta
  \[
    r\colon
    \begin{cases}
      x = 1 +  t \\
      y = 2 + 2t \\
      z = 1 +  t
    \end{cases}
  \]
  e o plano
  \[
    \alpha\colon x - 2y + 2z - 5 = 0
  \]
\end{question}

%------------------------------------------------------------------------------%
\begin{resolution}{Solução}
  Vetor diretor da reta
  \[
    \vec{v} = (1, 2, 1)
  \]

  \pause
  Vetor normal ao plano
  \[
    \vec{n} = (1, -2, 2)
  \]
\end{resolution}

%------------------------------------------------------------------------------%
\begin{resolution}{Solução}
  \begin{align*}
    \onslide<+->{\vec{v}\cdot\vec{n}}
    \onslide<+->{& = 1\cdot1+2\cdot(-2)+1\cdot2 \\}
    \onslide<+->{& = 1-4+2                      \\}
    \onslide<+->{& = -1}
  \end{align*}

  \begin{align*}
    \onslide<+->{\norm{\vec{v}}}
    \onslide<+->{& = \sqrt{1^2+2^2+1^2} \\}
    \onslide<+->{& = \sqrt{6}}
  \end{align*}

  \begin{align*}
    \onslide<+->{\norm{\vec{n}}}
    \onslide<+->{& = \sqrt{1^2+(-2)^2+2^2} \\}
    \onslide<+->{& = 3}
  \end{align*}
\end{resolution}

%------------------------------------------------------------------------------%
\begin{resolution}{Solução}
  \begin{align*}
    \onslide<+->{\sin\theta}
    \onslide<+->{& = \frac{|-1|} {\sqrt{6}\cdot3} \\}
    \onslide<+->{& = \frac{1}{3\sqrt{6}}}
  \end{align*}

  \[
    \onslide<+->{\theta = \arcsin\left(\frac{1}{3\sqrt{6}}\right)}
  \]
\end{resolution}

%------------------------------------------------------------------------------%
\endinput
